% Pursuing Stacks (cotes 134-2 to 134-8): archive cote 134-3, site page 93 (PDF page 94).
% notes for ch. II. Draft compared: ps2.tex:2190-2205 ; ps2.tex:2240-2257 ; ps2.tex:2290-2306 ; ps2.tex:2330-2347 ; ps2.tex:4400-4415.
% AI-transcribed by Claude Opus 5.5 (claude-opus-5-5, adaptive thinking, effort high), Message Batches API, batch ps1 (msgbatch_01Rv5hCaNTtUt4JHx1yqoHPL), 2026-11-12; image review recorded in records/CHECK_DECISIONS_134_3_4.tsv.
\scanpage{134-3/93}
II v

I am going to come back upon this in a later chapter in volume 2 of
Pursuing \over{S}{S}tacks (presumably chapter X : Back to topoï) - compare remarks
in section 35.

($^{33}$) We'll see later it \uline{is} stronger.

($^{34}$) It will become clear later \inserted{\deleted{(section 38, cor. 1)}} that it is'nt necessary indeed, \inserted{(Cf. end of section 39.} \deleted{For}
\deleted{instance, if A is any test category (for instance $\hand{\Delta}$) then so is $A_{/a}$}
\deleted{for any object in a , while obviously T 2 will no\over{\unsure{t}}{t} longer be true most}
\deleted{times (for instance if $A=\hand{\Delta}$, it is'nt except if $a = \hand{\Delta}_o$ , otherwise we}
\deleted{may find two subobjects of a, isomorphic to the final object $e=\hand{\Delta}_o$ , and}
\deleted{whose intersection (i.e. product in $(\hand{\Delta}_{/a})^{\hand{\wedge}}\,\hand{\approx}\hand{\Delta}^{\hand{\wedge}}_{/a}$) is the empty object ...}

(\hand{!}$^{35}$) This of course is \xout{false} coarsely false - the only types of \xout{\unsure{fini}}
limits (in the usual technical sense) which exist in (H\over[typed]{o}{o}t), as is well
known, are finite products and arbitrary sums. I was confusing with the
more sophisticated notion of ``homotopy limits'', which are altogether a
different story\hand{!}

($^{36}$) This is indeed easily checked, and will be done in part \over[typed]{\unsure{V}}{X}III
(volume 2 of Pursuing Stacks).

($^{37}$) (Added 9.4.1983) Presumably, this condition (i) is too weak as
stated, we'll have to add that $j_A=i_A^{*}$ is equally modelizing, cf. remark
3 \over{o}{i}\over{n}{n} \xout{page \unsure{1}} section 65 , D).

($^{38}$) Finally, it turns out later \inserted{(section 39)} that it \xout{was} \inserted[typed]{is} better to keep the notion
of a ``strict test category'' as we had it before (implying T 2), whereas
the notion of a``test category'' will be \deleted{however} slightly stronger than
before. For a\hand{n} \inserted{overall} review of the test notions (using the usual notion of weak
equivalence in the basic modelizer (Cat)\hand{),} see section 44 .

\mspage{\hand{93}}
